% !TEX root = ../print.tex
% Chapter 5, continued: section 5.3 of the article, invariant wedges (not in the draft). Provenance:
% the plan's item D3 (notes/PLAN-article-structure.md, section 5.3 and "The organizing idea"; queue
% Q-0266). Until Q-0278 (2026-10-04) this part was 02-wedges.tex, placed after the stable-cone nodes it
% reads; it now sits where the plan's outline puts it. Its last node, thm:wedge-axis, which \uses nodes of
% 05-orbits.tex and 05-smaller-groups.tex, moved to 5.7 (05-table.tex) on the author's decision of
% 2026-10-04. No label, statement or proof changed.
%
% Design. The order W of a wedge datum (def:gw-cascade-family) is, on PD(d), determined by the set of
% relative steps it allows. This part says which step sets are orders (prop:invariant-orders), proves
% the correspondence between closed invariant orders generated by cone fields and O(d)-invariant
% closed convex cones in Sym(d), equivalently permutation-symmetric cones in eigenvalue space
% (lem:spectral-cones, thm:cone-orders), answers the plan's question whether the stable step cone is
% the order generated by K_alpha (prop:stable-order-generated: yes, exactly, with no closure). The chain
% of wedges with the general linear row along it (thm:wedge-axis) is stated in 5.7. Everything here is
% proved in-house but one matrix inequality: since the author's decision of 2026-10-05 the
% logarithmic majorization of a product (lem:spectral-cones-products) is ledger entry A9 (Bhatia,
% Positive Definite Matrices, Ex. 6.5.3, p. 228, the weak form), and the nodes that read it through
% lem:spectral-cones(5) spend it; Ky Fan's minimum principle is proved, and prop:stable-order-generated
% spends nothing (its proof reads the stable orders through Phi_alpha and the Loewner order). The prior art (Lawson 1989; Hilgert-Hofmann-Lawson
% 1989; Hilgert-Olafsson 1997; Faraut-Koranyi 1994) is read at the page and recorded in
% records/affine/READING-*.md; no node's proof depends on it.

\section{Invariant wedges: orders on positive definite matrices and their cones}
\label{sec:invariant-wedges}

The order $W$ of a wedge datum (\ref{def:gw-cascade-family}) is invariant under a group containing the
dilations, and for the general linear group it is, on $\PD(d)$, determined by the set of relative steps
$M = \Sigma^{-1/2}\Sigma'\Sigma^{-1/2}$ it allows. This section asks which sets of steps give an order,
and how such an order is read infinitesimally. The steps do not compose by multiplication: the step of a
composite is $M^{1/2}NM^{1/2}$ up to a rotation of $N$, and its spectrum depends on how $M$ and $N$ sit
relative to each other (\ref{prop:invariant-orders}). An order generated by an $\Orth(d)$-invariant
cone field is exactly the order whose steps are the exponentials of the cone, and such cones are the
permutation-symmetric cones in eigenvalue space (\ref{lem:spectral-cones}, \ref{thm:cone-orders}). The
stable step cone of \ref{prop:stable-step-cone} is the order generated by the symmetric part of the
exponent cone $\mathfrak k_\alpha$, and so are the ray order and the Loewner order for their cones
(\ref{prop:stable-order-generated}). The cones of the module form a chain indexed by
$\alpha \in [0, \infty]$, and along it the families covariant under the general linear group shrink
from all self-decomposable orbits to the Gaussian, which is alone from $\alpha = 2$ on; that is
\ref{thm:wedge-axis}, stated in \ref{sec:landscape-table} once the families it reads are in place.

% not in the draft: the plan's item D3 (Q-0266). What writing it down found: the steps of a composite
% are not the product of the steps, and the spectrum of the composite is not a function of the spectra
% of the two steps (clause (3)'s example); transitivity is a condition on all relative positions, which
% is why the next node works with majorization rather than with spectra one by one.
\begin{proposition}[invariant relations on $\PD(d)$ and their steps]
\label{prop:invariant-orders}
\lean{AffineHemigroup.stepLE_iff_factor, AffineHemigroup.stepLE_invariant,
AffineHemigroup.invariant_relation_eq_stepLE, AffineHemigroup.stepLE_order_props,
AffineHemigroup.stepLE_compose, AffineHemigroup.stepLE_chain_of_mem,
AffineHemigroup.stepLE_trans_tfae, AffineHemigroup.step_composite_example,
AffineHemigroup.infCone_props}\leanok
Let $d \ge 1$ and write $\Sym(d)$ for the real symmetric $d \times d$ matrices. Call a set
$\mathcal K \subseteq \PD(d)$ a \emph{step set} if $QMQ^{\tp} \in \mathcal K$ whenever $M \in \mathcal K$
and $Q \in \Orth(d)$. For a step set write $\Sigma \le_{\mathcal K} \Sigma'$, for
$\Sigma, \Sigma' \in \PD(d)$, when the \emph{relative step} $\Sigma^{-1/2}\Sigma'\Sigma^{-1/2}$ lies in
$\mathcal K$, and put $\mathcal S_{\mathcal K} := \{A \in \GL(d) : AA^{\tp} \in \mathcal K\}$.
\begin{enumerate}
\item \emph{(Invariant relations are step sets.)} For $\Sigma, \Sigma' \in \PD(d)$ and every
$A \in \GL(d)$ with $AA^{\tp} = \Sigma$, $\Sigma \le_{\mathcal K} \Sigma'$ if and only if
$A^{-1}\Sigma'A^{-\tp} \in \mathcal K$. The relation $\le_{\mathcal K}$ is $\GL(d)$-invariant:
$\Sigma \le_{\mathcal K} \Sigma'$ if and only if $B\Sigma B^{\tp} \le_{\mathcal K} B\Sigma'B^{\tp}$, for
every $B \in \GL(d)$. Every $\GL(d)$-invariant relation $R$ on $\PD(d)$ is $\le_{\mathcal K}$ for
exactly one step set, $\mathcal K = \{M \in \PD(d) : I \,R\,M\}$.
\item \emph{(Order properties.)} $\le_{\mathcal K}$ is reflexive if and only if $I \in \mathcal K$; it is
antisymmetric if and only if no $M \ne I$ has both $M$ and $M^{-1}$ in $\mathcal K$; and it is closed in
$\PD(d) \times \PD(d)$ if and only if $\mathcal K$ is closed in $\PD(d)$.
\item \emph{(How steps compose.)} If $\Sigma \le_{\mathcal K} \Sigma' \le_{\mathcal K} \Sigma''$, $M$ is
the relative step of $(\Sigma, \Sigma')$ and $N$ that of $(\Sigma', \Sigma'')$, then the relative step
of $(\Sigma, \Sigma'')$ is $M^{1/2}N'M^{1/2}$ with $N' = Q^{\tp}NQ$ for some $Q \in \Orth(d)$, and its
eigenvalues are those of $MN'$. Conversely, for $M, N \in \mathcal K$ the triple $I$, $M$,
$M^{1/2}NM^{1/2}$ is such a chain. So $\le_{\mathcal K}$ is transitive if and only if
$M^{1/2}NM^{1/2} \in \mathcal K$ for all $M, N \in \mathcal K$, if and only if $\mathcal S_{\mathcal K}$
is closed under products. The spectrum of the composite is not determined by the spectra of the two
steps: for $d = 2$ the step $\diag(4,1)$ followed by $\diag(4,1)$ gives $\diag(16,1)$, and followed by
$\diag(1,4)$ gives $4I$.
\item \emph{(The infinitesimal cone.)} The set
$C(\mathcal K) := \{X \in \Sym(d) : e^{tX} \in \mathcal K \text{ for every } t \ge 0\}$ is invariant
under $X \mapsto QXQ^{\tp}$, $Q \in \Orth(d)$, and under $X \mapsto cX$, $c \ge 0$. It is closed when
$\mathcal K$ is closed in $\PD(d)$, and if $\le_{\mathcal K}$ is antisymmetric then
$X \in C(\mathcal K)$ and $-X \in C(\mathcal K)$ only for $X = 0$.
\end{enumerate}
\statusT\quad\emph{(Cites nothing beyond the spectral theorem for symmetric matrices. \textbf{What it is
worth.} It is the dictionary between the orders of the scheme and sets of matrices: a
$\GL(d)$-invariant order on $\PD(d) = \GL(d)/\Orth(d)$ is its set of relative steps at $I$, and the
order is transitive exactly when $\mathcal S_{\mathcal K}$ is a semigroup, which is the correspondence
between invariant orders on a homogeneous space $G/H$ and semigroups containing $H$ (in print:
\sscite{lawson1989ordered} Prop.~6.3, pp.~295--296; \sscite{hilgert1997causal} Lemma~2.3.1 and
Remark~2.3.2, pp.~43--44; read, not used). Clause (3) is the answer to ``the steps do not compose by
multiplication'': the composite is $M^{1/2}N'M^{1/2}$, symmetric, with the spectrum of $MN'$, and the
rotation $Q$ between the two eigenframes enters. So transitivity is a condition on all relative
positions of two steps: the composite's spectrum is not a function of the two spectra, so a set of
spectra closed under a fixed single-valued operation need not describe the order, and transitivity is
closure under the set-valued map taking two spectra to all spectra of $M^{1/2}QNQ^{\tp}M^{1/2}$,
$Q \in \Orth(d)$ (the multiplicative Horn set).
The next node handles this with majorization, which bounds the spectrum of the composite uniformly in
$Q$. The orders of \ref{thm:cascade-infinitely-divisible}(3)(c) are the $\le_{\mathcal K}$ with
$\mathcal K$ closed under powers $M^\tau$, $0 \le \tau \le 1$; the closed ones among them are the
orders of cone fields (\ref{thm:cone-orders}(2)). \textbf{Not claimed}: a description of the set of
spectra of $M^{1/2}QNQ^{\tp}M^{1/2}$ over all $Q$ for given spectra of $M$ and $N$ (a multiplicative
Horn problem; the module does not need it); and anything about relations on $\PSD(d)$ beyond
\ref{thm:cone-orders}(4).)}
\end{proposition}

% CHANGED Q-0281 wave 2: proved (AffineHemigroup.InvariantOrders; Lean core, guarded). Proof
% rewritten to the Lean route in (2): the reversed step through the factor Sigma^{1/2}M^{1/2} of
% Sigma' (by (1)) in place of the orthogonal conjugacy of isospectral matrices, and the closedness
% through the continuity of the square root on PSD(d), proved by compactness.
\begin{proof}
\leanok
\emph{(1).} Let $AA^{\tp} = \Sigma$ and $Q := \Sigma^{-1/2}A$; then $QQ^{\tp} =
\Sigma^{-1/2}AA^{\tp}\Sigma^{-1/2} = I$, so $Q \in \Orth(d)$ and $A^{-1}\Sigma'A^{-\tp} =
Q^{\tp}\Sigma^{-1/2}\Sigma'\Sigma^{-1/2}Q$, which lies in $\mathcal K$ exactly when the relative step does,
$\mathcal K$ being a step set. For $B \in \GL(d)$, $BA$ satisfies $(BA)(BA)^{\tp} = B\Sigma B^{\tp}$ and
$(BA)^{-1}B\Sigma'B^{\tp}(BA)^{-\tp} = A^{-1}\Sigma'A^{-\tp}$; with what was just shown, this is the
invariance. Let $R$ be a $\GL(d)$-invariant relation and $\mathcal K := \{M : I \,R\,M\}$. For
$Q \in \Orth(d)$, $QIQ^{\tp} = I$, so $I \,R\,M$ gives $I \,R\,QMQ^{\tp}$: $\mathcal K$ is a
step set. By invariance under $B = \Sigma^{-1/2}$, $\Sigma \,R\,\Sigma'$ if and only if
$I \,R\,\Sigma^{-1/2}\Sigma'\Sigma^{-1/2}$, that is, $\Sigma \le_{\mathcal K} \Sigma'$. A step set
is recovered from its relation as $\{M : I \le_{\mathcal K} M\}$, which gives uniqueness.

\emph{(2).} The relative step of $(\Sigma, \Sigma)$ is $I$. For antisymmetry let $M$ be the relative step
of $(\Sigma, \Sigma')$. The matrix $A := \Sigma^{1/2}M^{1/2}$ has $AA^{\tp} = \Sigma^{1/2}M\Sigma^{1/2} =
\Sigma'$ and $A^{-1}\Sigma A^{-\tp} = M^{-1/2}(\Sigma^{-1/2}\Sigma\Sigma^{-1/2})M^{-1/2} = M^{-1}$, so by
(1) $\Sigma' \le_{\mathcal K} \Sigma$ if and only if $M^{-1} \in \mathcal K$. The
relative step $M$ is $I$ exactly when $\Sigma = \Sigma'$, and every matrix of $\PD(d)$ is the relative step
of a pair, $N$ that of $(I, N)$. So $\le_{\mathcal K}$ is antisymmetric exactly when no $M \ne I$ has $M$
and $M^{-1}$ in $\mathcal K$. The map $\Sigma \mapsto \Sigma^{1/2}$ is continuous on $\PSD(d)$: if
$\Sigma_n \to \Sigma$, the entries of $\Sigma_n^{1/2}$ are bounded, $(\Sigma_n^{1/2})_{ij}^2 \le
(\Sigma_n)_{ii}$, and every limit point of $(\Sigma_n^{1/2})$ is positive semidefinite with square
$\Sigma$, hence is $\Sigma^{1/2}$ by the uniqueness of the positive square root. So
$(\Sigma, \Sigma') \mapsto \Sigma^{-1/2}\Sigma'\Sigma^{-1/2}$ is continuous on $\PD(d) \times \PD(d)$, and a closed
$\mathcal K$ gives a closed relation; conversely $\mathcal K$ is the slice $\{M : (I, M) \in {\le_{\mathcal K}}\}$
of the relation.

\emph{(3).} Put $A := \Sigma^{1/2}M^{1/2}$, so that $AA^{\tp} = \Sigma^{1/2}M\Sigma^{1/2} = \Sigma'$,
and, as in (1), $A = \Sigma'^{1/2}Q$ with $Q \in \Orth(d)$. Then $\Sigma'' = \Sigma'^{1/2}N\Sigma'^{1/2}
= AQ^{\tp}NQA^{\tp} = \Sigma^{1/2}M^{1/2}N'M^{1/2}\Sigma^{1/2}$, so the relative step of
$(\Sigma, \Sigma'')$ is $M^{1/2}N'M^{1/2} = M^{-1/2}(MN')M^{1/2}$, which has the eigenvalues of $MN'$.
The converse chain has relative steps $M$ and $M^{-1/2}(M^{1/2}NM^{1/2})M^{-1/2} = N$. So transitivity
is the stated closure property, $N' \in \mathcal K$ ranging over $\mathcal K$ as $N$ does. If it holds
and $A_1, A_2 \in \mathcal S_{\mathcal K}$, put $M_1 := A_1A_1^{\tp}$ and write $A_1 = M_1^{1/2}Q_1$ with
$Q_1 \in \Orth(d)$. Then $(A_1A_2)(A_1A_2)^{\tp} = M_1^{1/2}(Q_1A_2A_2^{\tp}Q_1^{\tp})M_1^{1/2}$ with
$Q_1A_2A_2^{\tp}Q_1^{\tp}$ in $\mathcal K$, so $A_1A_2 \in \mathcal S_{\mathcal K}$. Conversely, if $\mathcal S_{\mathcal K}$ is closed under
products and $M, N \in \mathcal K$, then $M^{1/2}, N^{1/2} \in \mathcal S_{\mathcal K}$, and so is
$M^{1/2}N^{1/2}$, that is, $M^{1/2}NM^{1/2} \in \mathcal K$. For the example,
$\diag(2,1)\diag(4,1)\diag(2,1) = \diag(16,1)$ and $\diag(2,1)\diag(1,4)\diag(2,1) = \diag(4,4)$.

\emph{(4).} $e^{tQXQ^{\tp}} = Qe^{tX}Q^{\tp}$ and $e^{t(cX)} = e^{(tc)X}$ give the two invariances. If
$X_n \to X$ with $X_n \in C(\mathcal K)$, then $e^{tX_n} \to e^{tX}$ in $\PD(d)$ for each $t$, so a closed
$\mathcal K$ contains $e^{tX}$. If $X$ and $-X$ lie in $C(\mathcal K)$, then $e^X$ and its inverse
$e^{-X}$ lie in $\mathcal K$, so antisymmetry gives $e^X = I$, hence $X = 0$, the exponential being
injective on $\Sym(d)$.
\end{proof}

% not in the draft: the plan's item D3 (Q-0266). The matrix inequalities behind the correspondence,
% proved rather than cited: Ky Fan's minimum principle for sums, and the multiplicative majorization of
% the spectrum of a product through compound matrices (Cauchy-Binet). Clause (3) is the dual description
% of a closed convex cone (Faraut-Koranyi Thm. I.1.1, p. 2, read and not used; the projection argument is
% written out). What writing it down found: a closed convex cone that is symmetric under permutations is
% automatically closed under majorization, so neither Rado's theorem nor the multiplicative Horn
% inequalities are needed.
% CHANGED Q-0281: typed (the statement skeleton). Clause (2), the products, is split to
% lem:spectral-cones-products (ADR-0025; records/formalization/PRICING.md: compound matrices and
% Cauchy-Binet are most of the node's cost); the text is unchanged and the \lean tag names the child's
% declaration too, and AffineHemigroup.eigUp_spec, which fixes the notation lambda-up.
% CHANGED Q-0281: blind review R3-1: \uses gains the split child lem:spectral-cones-products, which
% clause (5)'s proof reads ("(2) with (3)"); the child no longer \uses this node.
% CHANGED Q-0281 wave 2, group F: proved (AffineHemigroup.KyFan, AffineHemigroup.SpectralCones). The
% status moves to [A]: clause (2) is ledger A9 through the child (the author's decision of 2026-10-05),
% and clause (5)'s products read it; a declaration of spending, not a statement change. The proof is
% rewritten to the Lean route in (2) (A9 and a determinant), in (3) (a separating functional in
% place of the nearest point), in (4) (the cyclic shifts in place of all d! permutations) and in (5)
% (the continuity of ell through that of lambda-up).
\begin{lemma}[invariant cones, and the spectra of sums and products]
\label{lem:spectral-cones}
\uses{prop:invariant-orders,lem:spectral-cones-products,thm:cited-A9}
\lean{AffineHemigroup.eigUp_spec, AffineHemigroup.partialSum_eigUp_isLeast,
AffineHemigroup.partialSum_eigUp_lipschitz, AffineHemigroup.eigUp_add_majorized,
AffineHemigroup.logSpec_product_majorized, AffineHemigroup.symCone_mem_iff,
AffineHemigroup.symCone_majorization, AffineHemigroup.specCone_isClosedConvexCone,
AffineHemigroup.orthInvariant_cone_eq_specCone, AffineHemigroup.orthInvariant_pointed_cone_sign,
AffineHemigroup.expSet_specCone}\leanok
Let $d \ge 1$. For $X \in \Sym(d)$ write $\lambda^\uparrow(X) = (\lambda^\uparrow_1(X), \dots,
\lambda^\uparrow_d(X))$ for its eigenvalues in increasing order, and for $M \in \PD(d)$ put
$\ell(M) := \lambda^\uparrow(\log M) = (\log\lambda^\uparrow_1(M), \dots, \log\lambda^\uparrow_d(M))$. For
$x \in \RR^d$ let $x^\uparrow$ and $x^\downarrow$ be its increasing and decreasing rearrangements,
$s_k(x) := x^\uparrow_1 + \dots + x^\uparrow_k$, and write $x \prec y$ when $s_k(x) \ge s_k(y)$ for
$1 \le k < d$ and $s_d(x) = s_d(y)$. A set $D \subseteq \RR^d$ is \emph{symmetric} if it is invariant
under permutations of the coordinates, and $D^\# := \{w \in \RR^d : \langle w, x\rangle \ge 0 \text{ for
every } x \in D\}$.
\begin{enumerate}
\item \emph{(Sums.)} $s_k(\lambda^\uparrow(X)) = \min\{\tr(U^{\tp}XU) : U \in \RR^{d\times k},\
U^{\tp}U = I\}$. Consequently $|s_k(\lambda^\uparrow(X)) - s_k(\lambda^\uparrow(Y))| \le k\|X - Y\|$, with
$\|\cdot\|$ the operator norm, and $\lambda^\uparrow(X + Y) \prec \lambda^\uparrow(X) + \lambda^\uparrow(Y)$.
\item \emph{(Products.)} For $M, N \in \PD(d)$, $\ell(M^{1/2}NM^{1/2}) \prec \ell(M) + \ell(N)$.
\item \emph{(Symmetric cones are closed under majorization.)} Let $D \subseteq \RR^d$ be a symmetric
closed convex cone. Then $x \in D$ if and only if $\sum_iw^\downarrow_ix^\uparrow_i \ge 0$ for every
$w \in D^\#$. Consequently $y \in D$ and $x \prec y$ imply $x \in D$.
\item \emph{(Cones in $\Sym(d)$ and in eigenvalue space.)} For a symmetric closed convex cone
$D \subseteq \RR^d$, $C_D := \{X \in \Sym(d) : \lambda^\uparrow(X) \in D\}$ is a closed convex cone
invariant under $X \mapsto QXQ^{\tp}$, $Q \in \Orth(d)$. Every closed convex cone $C \subseteq \Sym(d)$
with that invariance is $C_D$ for exactly one such $D$, namely $D = \{x : \diag(x) \in C\}$, and
$C \cap (-C) = \{0\}$ if and only if $D \cap (-D) = \{0\}$. If moreover $C \ne \{0\}$ and
$C \cap (-C) = \{0\}$, then exactly one of $I$ and $-I$ lies in $C$, and $\tr X$ has that sign for
every $X \in C \setminus \{0\}$.
\item \emph{(Exponentials.)} For $C = C_D$ as in (4), $\exp C := \{e^X : X \in C\}$ equals
$\{M \in \PD(d) : \ell(M) \in D\}$; it is a step set (\ref{prop:invariant-orders}), closed in $\PD(d)$,
and $M^{1/2}NM^{1/2} \in \exp C$ whenever $M, N \in \exp C$.
\end{enumerate}
\statusA\quad\ledger{A9}\quad\emph{(\textbf{Assignment.} Clause (2) is
\ref{lem:spectral-cones-products}, which spends \ledger{A9}, the inequality for the eigenvalues of a
product (\sscite{bhatia2007positive} Ex.~6.5.3, (6.47), p.~228), in its weak form, for the partial sums of
the largest entries; its equality of totals is a determinant, proved. The last assertion of clause (5), the
products, reads clause (2) and so rests on A9. Everything else is $[T]$: clause (1), Ky Fan's minimum
principle, is proved here, and so are (3), (4) and the rest of (5). \textbf{What it is worth.} Clause (5)
is the transitivity of every order whose steps are the exponentials of an invariant convex cone; clause (4)
is the plan's ``permutation-symmetric cone in eigenvalue space''. Clauses (1) and (2) are classical
inequalities of matrix analysis (Ky Fan's minimum principle, and the majorization of the logarithmic
spectrum of a product); the first is proved here, the second is cited since the author's decision of
2026-10-05, its proof through compound matrices and the Cauchy--Binet formula costing more to formalize
than the rest of the section; the transitivity of (5) for cones with $C \cap (-C) = \{0\}$ is in
print in Lie-semigroup form (\ref{thm:cone-orders}'s annotation). The last sentence of (4) is the
in-house form of the fact that an invariant cone has a fixed point of the compact group
(\sscite{hilgert1997causal} Lemma~2.1.15, p.~36, read, not used): the only $\Orth(d)$-fixed directions in
$\Sym(d)$ are $\pm I$, so every such cone is built on the centre. \textbf{Not claimed}: the converse of
(3) for sets that are not cones; and a description of all spectra of $M^{1/2}NM^{1/2}$.)}
\end{lemma}

% CHANGED Q-0281 wave 2, group F: proof rewritten to the Lean route in (2), (3), (4) and (5).
\begin{proof}
\leanok
\emph{(1).} Write $X = \sum_j\lambda_jv_jv_j^{\tp}$ with $\lambda_j = \lambda^\uparrow_j(X)$ and $(v_j)$ an
orthonormal basis. For $U$ with $U^{\tp}U = I$, $\tr(U^{\tp}XU) = \sum_j\lambda_jp_j$ with
$p_j := |U^{\tp}v_j|^2$; since $UU^{\tp}$ is an orthogonal projection, $0 \le p_j \le 1$, and
$\sum_jp_j = \tr(U^{\tp}U) = k$. Then
\[
  \sum_j\lambda_jp_j - \sum_{j\le k}\lambda_j = \sum_{j\le k}\lambda_j(p_j - 1) + \sum_{j>k}\lambda_jp_j
  \ge \lambda_k\sum_{j\le k}(p_j - 1) + \lambda_k\sum_{j>k}p_j = 0 ,
\]
using $p_j - 1 \le 0$ and $\lambda_j \le \lambda_k$ for $j \le k$, and $\lambda_j \ge \lambda_k$ for
$j > k$; equality holds for $U = (v_1, \dots, v_k)$. Each $X \mapsto \tr(U^{\tp}XU)$ is linear and
$k\|\cdot\|$-Lipschitz, so their minimum is $k\|\cdot\|$-Lipschitz, and the minimum of a sum is at least
the sum of the minima: $s_k(\lambda^\uparrow(X+Y)) \ge s_k(\lambda^\uparrow(X)) +
s_k(\lambda^\uparrow(Y))$, with equality at $k = d$, where both sides are traces. The vector
$\lambda^\uparrow(X) + \lambda^\uparrow(Y)$ is increasing, so its $s_k$ is the right-hand side.

\emph{(2).} Let $P := M^{1/2}NM^{1/2}$. The products $MN = M^{1/2}(M^{1/2}N)$ and $(M^{1/2}N)M^{1/2} = P$
have the same characteristic polynomial, so the eigenvalues of $MN$, with multiplicity, are those of $P$.
By the inequality for the eigenvalues of a product in its weak form (\ledger{A9}), for $1 \le k \le d$
\[
  \sum_{i\le k}\log\lambda^\downarrow_i(P) \le \sum_{i\le k}\bigl(\log\lambda^\downarrow_i(M) +
  \log\lambda^\downarrow_i(N)\bigr) ,
\]
with $\lambda^\downarrow$ the eigenvalues in decreasing order: the sum of the $k$ largest entries of
$\ell(P)$ is at most that of $\ell(M) + \ell(N)$, both vectors being increasing. The totals are equal,
$\sum_i\ell_i(P) = \log\det P = \log(\det M\det N) = \sum_i(\ell_i(M) + \ell_i(N))$. Subtracting from the
totals, $s_j(\ell(P)) \ge s_j(\ell(M) + \ell(N))$ for $j = d - k$, $1 \le j < d$, with equality at $j = d$,
which is the claim.

\emph{(3).} \emph{The dual description.} If $x \in D$ and $w \in D^\#$, every rearrangement of $x$ lies in
$D$, so $\langle w, \sigma x\rangle \ge 0$ for every permutation $\sigma$; the number
$\sum_iw^\downarrow_ix^\uparrow_i$ is one of these, the largest entries of $w$ meeting the smallest of $x$,
and it is the least of them (the rearrangement inequality). Conversely let $x \notin D$. As $D$ is closed
and convex, a linear functional $f$ separates $x$ from $D$ strictly: $f(x) < u < f(z)$ for a number $u$ and
every $z \in D$ (Hahn--Banach). Write $f(z) = \langle w, z\rangle$. With $z = 0$, $u < 0$; and if
$\langle w, z\rangle < 0$ for some $z \in D$, then $cz \in D$ with $c := u/\langle w, z\rangle > 0$ gives
$u < \langle w, cz\rangle = u$, which is absurd. So $w \in D^\#$ and $\langle w, x\rangle < u < 0$, and by
the rearrangement inequality $\sum_iw^\downarrow_ix^\uparrow_i \le \langle w, x\rangle < 0$.
\emph{Majorization.} By Abel summation, for $w \in \RR^d$ and $x \in \RR^d$,
\[
  \sum_{i=1}^dw^\downarrow_ix^\uparrow_i = \sum_{k=1}^{d-1}\bigl(w^\downarrow_k - w^\downarrow_{k+1}\bigr)s_k(x)
  + w^\downarrow_d\,s_d(x) ,
\]
with every coefficient $w^\downarrow_k - w^\downarrow_{k+1} \ge 0$. If $x \prec y$, each term for $x$ is at
least the corresponding term for $y$ (the last ones are equal), so
$\sum_iw^\downarrow_ix^\uparrow_i \ge \sum_iw^\downarrow_iy^\uparrow_i \ge 0$ for $w \in D^\#$ when
$y \in D$, and $x \in D$ by the dual description.

\emph{(4).} $\lambda^\uparrow(QXQ^{\tp}) = \lambda^\uparrow(X)$ and $\lambda^\uparrow(cX) = c\lambda^\uparrow(X)$
for $c \ge 0$. Each $\lambda^\uparrow_k = s_k - s_{k-1}$ is continuous by (1), and $D$ is closed, so $C_D$
is closed. If $X, Y \in C_D$, then $\lambda^\uparrow(X) + \lambda^\uparrow(Y) \in D$, $D$ being a convex
cone, and (1) with (3) gives $\lambda^\uparrow(X+Y) \in D$. So $C_D$ is a convex cone. Given $C$, the set
$D := \{x : \diag(x) \in C\}$ is a closed convex cone, symmetric because permutation matrices are
orthogonal. Writing $X = Q\diag(\lambda^\uparrow(X))Q^{\tp}$, $X \in C$ if and only if $\lambda^\uparrow(X)
\in D$, so $C = C_D$. And $D$ is determined by $C_D$ as $\{x : \diag(x) \in C_D\}$. If $0 \ne x \in D \cap
(-D)$ then $\diag(x) \in C \cap (-C)$; if $0 \ne X \in C \cap (-C)$ then $\lambda^\uparrow(X)$ and
$\lambda^\uparrow(-X)$, a rearrangement of $-\lambda^\uparrow(X)$, lie in $D$, so $0 \ne \lambda^\uparrow(X)
\in D \cap (-D)$. For the last sentence let $0 \ne X \in C$ and $x := \lambda^\uparrow(X) \ne 0$. For
$m = 0, \dots, d-1$ let $\sigma_mx$ be the cyclic shift $(\sigma_mx)_i := x_{i+m}$, indices modulo $d$.
Each $\sigma_mx$ lies in $D$, and their sum is $(\tr X)\mathbf 1$, with $\mathbf 1 = (1, \dots, 1)$, since
every coordinate of it is $\sum_jx_j$; so $(\tr X)\mathbf 1 \in D$. If $\tr X = 0$ this sum is $0$, so
$x = \sigma_0x = -\sum_{m\ne0}\sigma_mx \in -D$ and $x \in D \cap (-D)$, which is excluded. So $\tr X \ne 0$,
and $(\operatorname{sign}\tr X)\mathbf 1 \in D$, that is, $(\operatorname{sign}\tr X)\,I \in C$. As $I$ and
$-I$ cannot both lie in $C$, the sign is the same for every nonzero $X \in C$, which is the claim.

\emph{(5).} For $X \in \Sym(d)$, $e^X$ has the eigenvalues $e^{\lambda_i(X)}$, and $\log$ is the inverse of
$\exp$ from $\Sym(d)$ onto $\PD(d)$; so $M \in \exp C$ if and only if $\log M \in C$, that is,
$\ell(M) \in D$. Orthogonal conjugation commutes with $\exp$, so $\exp C$ is a step set. On $\PD(d)$,
$\ell(M) = (\log\lambda^\uparrow_i(M))_i$ is continuous, each $\lambda^\uparrow_i$ being continuous by (1)
and positive there, and $D$ is closed, so $\exp C = \{M : \ell(M) \in D\}$ is closed in $\PD(d)$. For
$M, N \in \exp C$, $\ell(M) + \ell(N) \in D$, and (2) with (3) gives $\ell(M^{1/2}NM^{1/2}) \in D$.
\end{proof}

% CHANGED Q-0281: split from lem:spectral-cones (ADR-0025; records/formalization/PRICING.md): clause
% (2), copied verbatim. Its proof through compound matrices and Cauchy-Binet is most of the parent's
% cost, and Mathlib has neither; the parent's text and numbering are unchanged.
% CHANGED Q-0281: blind review R3-1: the edge runs parent -> child; this node reads only the parent's
% notation (ell, the majorization), never a claim, and the parent's clause (5) reads this node.
% CHANGED Q-0281 wave 2, group F: proved (AffineHemigroup.SpectralProducts) on ledger A9, the
% author's decision of 2026-10-05 in place of the compound matrices; the status moves to [A] with its
% assignment, a declaration of spending, not a statement change.
\begin{lemma}[the logarithmic spectrum of a product]
\label{lem:spectral-cones-products}
\uses{thm:cited-A9}
\lean{AffineHemigroup.logSpec_product_majorized}\leanok
Let $d \ge 1$, and use the notation of \ref{lem:spectral-cones}. For $M, N \in \PD(d)$,
$\ell(M^{1/2}NM^{1/2}) \prec \ell(M) + \ell(N)$.
\statusA\quad\ledger{A9}\quad\emph{(\textbf{Assignment.} \ledger{A9} carries the partial sums: for
$1 \le k \le d$, the sum of the $k$ largest entries of $\log\lambda(MN)$ is at most that of
$\log\lambda^\downarrow(M) + \log\lambda^\downarrow(N)$, which is the second half of
\sscite{bhatia2007positive} Ex.~6.5.3, (6.47), p.~228, in its weak form; the source leaves its symbol
undefined, and either reading implies these inequalities. It does not carry the equality of the totals,
which is $\det(M^{1/2}NM^{1/2}) = \det M\det N$ and is proved, nor the passage from the eigenvalues of
$MN$ to those of $M^{1/2}NM^{1/2}$, nor that from the largest entries to the module's $s_k$ of the
smallest, both proved. Clause (2) of \ref{lem:spectral-cones}, split off on 2026-10-04 (Q-0281,
ADR-0025); the numbering stays the parent's.)}
\end{lemma}

\begin{proof}
\leanok
This is the part of the proof of \ref{lem:spectral-cones} headed (2).
\end{proof}

% CHANGED Q-0281 wave 2: blind review W2H-2 (a delegated split, ADR-0025). The definitions of the cone
% field, the C-causal curves and the orders preceq_C and preceq^+_C, copied verbatim from
% thm:cone-orders (its preamble and the second sentence of its (4)), so that the nodes whose statements
% use this vocabulary (thm:cone-orders, thm:wedge-axis, lem:wedge-axis-chain) \uses a definition and not
% an [A] theorem. The parent's text and numbering are unchanged: it restates the definitions, as
% lem:cone-orders-psd(2) restates preceq^+_C, where they claim it is a partial order. The definition
% of preceq^+_C needs no hypothesis beyond the standing one, and is stated without I in C and
% C cap (-C) = {0}, as the Lean's causalPlusLE is. The cone field has no declaration of its own: the
% Lean reads it inline, as the image of C under X -> Sigma^{1/2} X Sigma^{1/2} (coneField_eq, in
% thm:cone-orders, which also carries the claim that it equals A C A^T). The tagged declarations are
% definitions of the proved library, sorry-free, so the node is \leanok.
\begin{definition}[cone fields, causal curves and their orders]
\label{def:cone-orders-curves}
\lean{AffineHemigroup.mderiv, AffineHemigroup.IsCausalCurve, AffineHemigroup.causalLE,
AffineHemigroup.causalPlusLE}\leanok
Let $d \ge 1$ and let $C \subseteq \Sym(d)$ be a closed convex cone invariant under $X \mapsto QXQ^{\tp}$,
$Q \in \Orth(d)$. Its \emph{cone field} assigns to $\Sigma \in \PD(d)$ the cone
$\Sigma^{1/2}C\Sigma^{1/2}$ of tangent vectors. An absolutely continuous $\gamma : [0,1] \to \PD(d)$ is
\emph{$C$-causal} if $\gamma(t)^{-1/2}\gamma'(t)\gamma(t)^{-1/2} \in C$ for almost every $t$; write
$\Sigma \preceq_C \Sigma'$ when a $C$-causal curve runs from $\Sigma$ to $\Sigma'$. For
$\Sigma, \Sigma' \in \PSD(d)$ let $\Sigma \preceq^+_C \Sigma'$ when $\Sigma = 0$, or $\Sigma' = c\,\Sigma$ for
some $c \ge 1$, or $\Sigma, \Sigma' \in \PD(d)$ and $\Sigma \preceq_C \Sigma'$.
\statusT\quad\emph{(Vocabulary, and no claim. The curves are read entrywise: $\gamma$ is absolutely
continuous when its entries are, and $\gamma'$ is the matrix of the derivatives of the entries, which exist
almost everywhere. What kind of relations $\preceq_C$ and $\preceq^+_C$ are is the content of
\ref{thm:cone-orders}(1) and (4), which spend entry A9; that $\preceq^+_C$ is a partial order when $I \in C$
and $C \cap (-C) = \{0\}$ is also proved without it, in the proof of \ref{thm:wedge-axis}. Split from \ref{thm:cone-orders} on 2026-10-06 (Q-0281, ADR-0025).)}
\end{definition}

% not in the draft: the plan's item D3 (Q-0266), "the correspondence orders <-> O(d)-invariant cones (as
% far as it is true)". What writing it down found: the correspondence is exact between invariant cones
% and the closed orders whose step sets are closed under powers (clause (2)); a closed invariant order in
% general contains the order of its cone and need not equal it (clause (3)); the order generated by a cone
% field needs no closure here, because the monotone functionals of the proof of (1) cut it out.
% CHANGED Q-0281: typed (the statement skeleton). Clause (2) is split to lem:cone-orders-closed and
% clauses (3)-(4) to lem:cone-orders-psd (ADR-0025; records/formalization/PRICING.md: (2) needs the
% Lie product formula, which Mathlib lacks; (3)-(4) are of a different kind); the text is unchanged
% and the \lean tag names the children's declarations too. C-causal curves are read entrywise
% (records/formalization/SKELETON.md, dispatch 1, reading 13).
% CHANGED Q-0281: \uses re-pointed to the split child it reads (thm:cascade-infinitely-divisible (3)(b)-(c) beside (2); lem:spectral-cones(2) beside the others).
% CHANGED Q-0281 wave 2, group F: proved (AffineHemigroup.ConeOrders, .LieProduct, .ConeOrdersClosed,
% .ConeOrdersPSD). The status moves to [A]: clause (1)'s superadditivity of F_w reads
% lem:spectral-cones(2), ledger A9 since the author's decision of 2026-10-05, and (2) and (4) read (1);
% a declaration of spending, not a statement change. The proof is rewritten to the Lean route in (1)
% (the geodesic as AA^T; the absolute continuity of g and its slopes through lambda-up and the scalar
% log), in (2) (the Lie product formula from a derivative), in (3) (why no cone generates the order)
% and in (4) (the chains through lem:path-condition-orders(c)).
% CHANGED Q-0281 wave 2: blind review W2F-5: clause (3)'s parenthetical names the geodesic of (1), which
% the proof uses, not (2), which spends A9; the statement narrows to the route the proof takes.
% CHANGED Q-0281 wave 2: blind review W2H-2: \uses gains def:cone-orders-curves, the split-off
% definitions of the cone field, the C-causal curves, preceq_C and preceq^+_C; the statement text,
% shared with the paper, is unchanged and restates them.
% CHANGED Q-0469 (wording only; records/affine/GATE-affine.md, Q-0317 rows "§4 21" and "§4 22"):
% (i) "Use the notation of prop:invariant-orders and lem:spectral-cones" names what it takes from each:
% the step-set vocabulary of the first, and exp C and lambda-up_1 (the least eigenvalue), which the
% statement writes out, from the second, which the paper states only in an appendix. (ii) Clause (3)'s
% parenthesis, "(by the geodesic of (1), C' would lie in C(K) = {0}, and K != {I})", a proof inside the
% statement, goes; the proof of (3) carries that argument. The same parenthesis goes from the verbatim
% copy in lem:cone-orders-psd(1). No claim changes; \uses unchanged (the proof reads lem:spectral-cones).
\begin{theorem}[invariant orders generated by cone fields]
\label{thm:cone-orders}
\uses{def:cone-orders-curves,prop:invariant-orders,lem:spectral-cones,def:gw-cascade-family,thm:cascade-infinitely-divisible,lem:path-condition-orders,lem:spectral-cones-products,thm:cited-A9}
\lean{AffineHemigroup.coneField_eq, AffineHemigroup.causalLE_iff_expSet,
AffineHemigroup.causalLE_order, AffineHemigroup.closed_order_tangent_wedge,
AffineHemigroup.closed_order_cone_tfae, AffineHemigroup.closed_order_without_cone,
AffineHemigroup.causalPlus_wedge, AffineHemigroup.causalPlus_ray}\leanok
Let $d \ge 1$ and let $C \subseteq \Sym(d)$ be a closed convex cone invariant under $X \mapsto QXQ^{\tp}$,
$Q \in \Orth(d)$. Its \emph{cone field} assigns to $\Sigma \in \PD(d)$ the cone
$\Sigma^{1/2}C\Sigma^{1/2}$ of tangent vectors, which equals $ACA^{\tp}$ for every $A$ with
$AA^{\tp} = \Sigma$, so that $B(\Sigma^{1/2}C\Sigma^{1/2})B^{\tp}$ is the cone at $B\Sigma B^{\tp}$ for
$B \in \GL(d)$. An absolutely continuous $\gamma : [0,1] \to \PD(d)$ is \emph{$C$-causal} if
$\gamma(t)^{-1/2}\gamma'(t)\gamma(t)^{-1/2} \in C$ for almost every $t$; write $\Sigma \preceq_C \Sigma'$
when a $C$-causal curve runs from $\Sigma$ to $\Sigma'$. Use the notation of \ref{prop:invariant-orders}
(step sets $\mathcal K$, their relations $\le_{\mathcal K}$, the sets $\mathcal S_{\mathcal K}$ and the cones
$C(\mathcal K)$), write $\exp C := \{e^X : X \in C\}$, and write $\lambda^\uparrow_1(M)$ for the least
eigenvalue of $M$.
\begin{enumerate}
\item \emph{(The order of a cone field.)} $\Sigma \preceq_C \Sigma'$ if and only if the relative step
$M = \Sigma^{-1/2}\Sigma'\Sigma^{-1/2}$ lies in $\exp C$, and then $\gamma(t) := \Sigma^{1/2}M^t\Sigma^{1/2}$
is $C$-causal. So $\preceq_C$ is $\le_{\exp C}$: a closed $\GL(d)$-invariant preorder, which is a partial
order if and only if $C \cap (-C) = \{0\}$; and $C(\exp C) = C$.
\item \emph{(Closed invariant orders.)} Let $\mathcal K$ be a closed step set for which $\le_{\mathcal K}$
is a partial order. Then
$\{B \in \gl(d) : e^{tB} \in \mathcal S_{\mathcal K} \text{ for every } t \ge 0\} =
\{B : \tfrac12(B + B^{\tp}) \in C(\mathcal K)\}$, and this is a closed convex cone containing every
skew-symmetric matrix; $C(\mathcal K)$ is an $\Orth(d)$-invariant closed convex cone with
$C(\mathcal K) \cap (-C(\mathcal K)) = \{0\}$; and $\preceq_{C(\mathcal K)}$ is contained in
$\le_{\mathcal K}$. The following are equivalent: (a) $\le_{\mathcal K}$ is $\preceq_{C'}$ for some cone
$C'$ as above; (b) $\mathcal K = \exp C(\mathcal K)$; (c) $M^\tau \in \mathcal K$ whenever
$M \in \mathcal K$ and $0 \le \tau \le 1$. When they hold, $C' = C(\mathcal K)$.
\item \emph{(A closed order without a cone.)} $\mathcal K := \{I\} \cup \{M \in \PD(d) :
\lambda^\uparrow_1(M) \ge 2\}$ is a closed step set, $\le_{\mathcal K}$ is a partial order, and
$C(\mathcal K) = \{0\}$, so that $\preceq_{C(\mathcal K)}$ is equality; hence $\le_{\mathcal K}$ is
$\preceq_{C'}$ for no cone $C'$.
\item \emph{(On $\PSD(d)$.)} Suppose $I \in C$ and $C \cap (-C) = \{0\}$. For $\Sigma, \Sigma' \in
\PSD(d)$ let $\Sigma \preceq^+_C \Sigma'$ when $\Sigma = 0$, or $\Sigma' = c\,\Sigma$ for some $c \ge 1$,
or $\Sigma, \Sigma' \in \PD(d)$ and $\Sigma \preceq_C \Sigma'$. Then $\preceq^+_C$ is a $\GL(d)$-invariant
partial order on $\PSD(d)$ that contains the rays, every pair $\Sigma \preceq^+_C \Sigma'$ is joined by a
chain as in \ref{thm:cascade-infinitely-divisible}(2), and $(G, \preceq^+_C)$ is a wedge datum with the
path condition (P) for every closed subgroup $G$ of $\GL(d)$ containing the dilations. It is contained in
the Loewner order if and only if $C \subseteq \PSD(d)$. For $C = \{cI : c \ge 0\}$ it is the ray order of
\ref{def:gw-cascade-family}.
\end{enumerate}
\statusA\quad\ledger{A9}\quad\emph{(\textbf{Assignment.} \ledger{A9}, the inequality for the
eigenvalues of a product in its weak form (\sscite{bhatia2007positive} Ex.~6.5.3, (6.47), p.~228), carries the
superadditivity of the functionals $F_w$ in the proof of (1), through \ref{lem:spectral-cones}(2), and
through (1) the transitivity of $\le_{\exp C}$, clause (2)'s inclusion and equivalence, and clause (4).
Everything else is $[T]$: the cone field, the geodesic of (1), the absolute continuity and the slopes of
the functionals along a causal curve, the Lie product formula (proved below), and clause (3), which reads
only the geodesic. \textbf{What it is worth.} It is the
correspondence the plan asks for, as far as it is true. An $\Orth(d)$-invariant closed convex cone in
$\Sym(d)$, equivalently a permutation-symmetric one in eigenvalue space (\ref{lem:spectral-cones}(4)),
generates an invariant order whose steps are exactly the exponentials of the cone, reached along the
geodesics $\Sigma^{1/2}M^t\Sigma^{1/2}$ and needing no closure; and a closed invariant order comes from a
cone exactly when its steps are closed under fractional powers, which is the class of closed
orders among those of \ref{thm:cascade-infinitely-divisible}(3)(c). Clause (3) shows that the restriction is needed: a closed
invariant order can have no infinitesimal part at all. The set in (2) is the tangent wedge
$\mathbf L(\mathcal S_{\mathcal K})$ of the closed semigroup $\mathcal S_{\mathcal K}$ in the sense of
\sscite{hilgert1989lie} Prop.~V.1.7, pp.~375--376. Clause (4) is the extension of
\ref{prop:similarity-stable-order-families}(1) to every cone; at $C = \PSD(d)$ it adds to the Loewner
order of $\PD(d)$ only the origin and the rays, so it is smaller than the Loewner order of $\PSD(d)$,
which has steps from nonzero singular matrices to definite ones. \textbf{In print.} The general frame is
Lawson's (\sscite{lawson1989ordered}): invariant cone fields on $G/H$ correspond to $\mathrm{Ad}(H)$-invariant
Lie wedges with edge the Lie algebra of $H$ (Thm.~4.9, p.~290), the conal order of a left-invariant wedge
field lies between the semigroup generated by $\exp W$ and its closure (Prop.~6.4, p.~296), and the
infinitesimally generated closed invariant orders correspond to the global Lie wedges (Thms.~7.4--7.5,
pp.~300--302, for $H$ closed and connected, here $\mathrm{SO}(d)$ in $\GL(d)^+$, whose quotient is again
$\PD(d)$). Hilgert--Ólafsson state the bijection between $H$-invariant regular cones and invariant regular
causal structures (\sscite{hilgert1997causal} Thm.~2.2.2, p.~40, with absolutely continuous causal curves on
p.~41, where they note that the strict causal relation need not be closed). For pointed cones, the
transitivity in (1) is \sscite{hilgert1989lie} Cor.~V.4.58, p.~440 (with Thm.~V.4.57, pp.~438--440): for a
Lie group with involution $\sigma$, fixed group $H$, and a pointed $\mathrm{Ad}(H_0)$-invariant cone $C$ in the
$(-1)$-eigenspace whose elements have real $\mathrm{ad}$-spectrum, $S = (\exp C)H_0$ is a semigroup with
tangent wedge $C \oplus \mathbf L(H)$ provided $S$ is closed. Read with $\sigma(A) = A^{-\tp}$ on $\GL(d)$,
$H = \Orth(d)$, $H_0 = \mathrm{SO}(d)$, its hypotheses hold and $S = (\exp C)\mathrm{SO}(d)$ is closed by the
polar decomposition, so it gives the transitivity in (1) for pointed cones; with Lawson's Props.~6.4--6.5
it would give the rest of (1) for piecewise smooth curves, a deduction not carried out here. For wedges, pointed
or not, \sscite{lawson1994polar} Thm.~3.1, pp.~198--199, derives the closedness: with $\mathrm{ad}\,X$ of real
spectrum on the wedge (Lem.~4.1 and Example~4.2, p.~203, for symmetric $X$) and the centre of $\gl(d)$ being
$\RR I$, which gives his condition (4), $(\exp C)\Orth(d) = \mathcal S_{\exp C}$ is a closed semigroup with subtangent
wedge $C \oplus \mathfrak{so}(d)$, which by \ref{prop:invariant-orders}(3) is the transitivity
of $\le_{\exp C}$ for every cone here; that paper has no order on the quotient, and its route is the Bony--Brezis
theorem through Cor.~V.4.58. Global wedges are
characterized by exact positive $1$-forms in \sscite{hilgert1989lie} Thm.~VI.5.1, p.~522; and the abelian example $\{(0,0)\} \cup {]1,\infty[^2}$ of
\sscite{hilgert1989lie} p.~376, a semigroup with tangent wedge $\{0\}$, is the model of (3). None of this is
used. \textbf{New here}: the reading on $\PD(d)$ through step sets and their spectra, the proof by
majorization (which needs neither the invariance theorem for wedges nor closedness as a hypothesis), the
equality of the order with $\exp C$ on $\PD(d)$ (the absolutely continuous curves
themselves are Hilgert--Ólafsson's), and the
characterization (2)(c). \textbf{Not claimed}: orders that are not closed; cone fields that are not
$\GL(d)$-invariant; relations on $\PSD(d)$ other than (4).)}
\end{theorem}

% CHANGED Q-0281 wave 2, group F: proof rewritten to the Lean route in (1), (2), (3) and (4).
\begin{proof}
\leanok
Write $C = C_D$ (\ref{lem:spectral-cones}(4)) and, for $w \in D^\#$ and $\Sigma, \Sigma' \in \PD(d)$,
\[
  F_w(\Sigma, \Sigma') := \sum_{i=1}^dw^\downarrow_i\,\ell_i\bigl(\Sigma^{-1/2}\Sigma'\Sigma^{-1/2}\bigr)
  = \sum_{k=1}^{d-1}\bigl(w^\downarrow_k - w^\downarrow_{k+1}\bigr)s_k(\ell) + w^\downarrow_d\,s_d(\ell) ,
\]
$\ell$ the logarithmic spectrum of the relative step, increasing. By \ref{lem:spectral-cones}(5) and (3),
$\Sigma \le_{\exp C} \Sigma'$ if and only if $F_w(\Sigma, \Sigma') \ge 0$ for every $w \in D^\#$.
\emph{Superadditivity.} For any $\Sigma, \Sigma', \Sigma'' \in \PD(d)$, with $M$, $N$ the relative steps of
$(\Sigma, \Sigma')$, $(\Sigma', \Sigma'')$, the relative step of $(\Sigma, \Sigma'')$ is
$M^{1/2}N'M^{1/2}$ with $N'$ orthogonally conjugate to $N$ (the computation of
\ref{prop:invariant-orders}(3), which does not use the order); by \ref{lem:spectral-cones}(2) its $s_k$ is
at least $s_k(\ell(M)) + s_k(\ell(N))$ for $k < d$, with equality at $k = d$. The coefficients
$w^\downarrow_k - w^\downarrow_{k+1}$ are nonnegative, so
$F_w(\Sigma, \Sigma'') \ge F_w(\Sigma, \Sigma') + F_w(\Sigma', \Sigma'')$.

\emph{(1), the geodesic.} Let $M \in \exp C$, $X := \log M \in C$ and $\gamma(t) = \Sigma^{1/2}e^{tX}\Sigma^{1/2}$.
With $A := \Sigma^{1/2}e^{tX/2}$, $\gamma(t) = AA^{\tp}$ and $\gamma'(t) = \Sigma^{1/2}Xe^{tX}\Sigma^{1/2} =
AXA^{\tp}$, so $\gamma(t)^{-1/2}\gamma'(t)\gamma(t)^{-1/2} = QXQ^{\tp}$ with $Q := \gamma(t)^{-1/2}A$, and
$QQ^{\tp} = \gamma(t)^{-1/2}AA^{\tp}\gamma(t)^{-1/2} = I$; so it lies in $C$. The entries of $\gamma$ are, in an
eigenbasis of $X$, finite sums of exponentials $e^{tx_k}$, so $\gamma$ is smooth.
\emph{(1), causal curves end in the order.} Let $\gamma$ be $C$-causal and $w \in D^\#$, and put
$g(t) := F_w(\gamma(0), \gamma(t))$, so $g(0) = 0$. On $[0,1]$, $g(t) = \sum_iw^\downarrow_i\log
\lambda^\uparrow_i(N(t))$ with $N(t) := \gamma(0)^{-1/2}\gamma(t)\gamma(0)^{-1/2}$. The least eigenvalue of
$N(t)$ is continuous in $t$ (\ref{lem:spectral-cones}(1)) and positive, so it is at least some $m > 0$ on
$[0,1]$, and $\log$ is $(1/m)$-Lipschitz on $[m, \infty)$. Each $\lambda^\uparrow_i$ is Lipschitz in the
operator norm by \ref{lem:spectral-cones}(1), and the operator norm of $N(s) - N(t)$ is at most
$\|\gamma(0)^{-1/2}\|^2\sum_{a,b}|\gamma_{ab}(s) - \gamma_{ab}(t)|$. So $|g(s) - g(t)|$ is bounded by a
constant times $\sum_{a,b}|\gamma_{ab}(s) - \gamma_{ab}(t)|$, and $g$ is absolutely continuous because the
entries of $\gamma$ are. Let $t \in (0,1)$ be a point at which the entries of $\gamma$ are differentiable and
$Y := \gamma(t)^{-1/2}\gamma'(t)\gamma(t)^{-1/2} \in C$; almost every $t$ is such. For $h > 0$,
$N_h := \gamma(t)^{-1/2}\gamma(t+h)\gamma(t)^{-1/2} = I + hY + o(h)$, and as $\lambda^\uparrow(I + hY) =
\mathbf 1 + h\lambda^\uparrow(Y)$, the Lipschitz bound of \ref{lem:spectral-cones}(1) gives
$x_i(h) := \lambda^\uparrow_i(N_h) = 1 + h\lambda^\uparrow_i(Y) + o(h)$. From $(x - 1)/x \le \log x \le x - 1$
for $x > 0$, $\log x_i(h) = h\lambda^\uparrow_i(Y) + o(h)$. Superadditivity gives
$g(t+h) - g(t) \ge F_w(\gamma(t), \gamma(t+h))$, hence, where $g$ is differentiable at $t$,
\[
  g'(t) \ge \lim_{h\downarrow0}\frac1hF_w(\gamma(t), \gamma(t+h)) = \sum_iw^\downarrow_i\lambda^\uparrow_i(Y) \ge 0 ,
\]
by \ref{lem:spectral-cones}(3), as $\lambda^\uparrow(Y) \in D$. An absolutely continuous function is the
integral of its derivative, which exists almost everywhere; so $g(1) = \int_0^1g' \ge 0$, and as $w$ was
arbitrary, $\gamma(0) \le_{\exp C} \gamma(1)$. \emph{(1), the rest.} $\exp C$ is a closed step set
containing $I$ (\ref{lem:spectral-cones}(5)), so $\le_{\exp C}$ is closed, reflexive and, by
\ref{lem:spectral-cones}(5) and \ref{prop:invariant-orders}(3), transitive; $M$ and $M^{-1}$ lie in
$\exp C$ exactly when $\log M \in C \cap (-C)$, which with \ref{prop:invariant-orders}(2) is the criterion
for antisymmetry; and $e^{tX} \in \exp C$ if and only if $tX \in C$, $\log$ being injective, so
$C(\exp C) = C$.

\emph{(2).} Let $\mathcal L$ be the set of $B$ named. $\mathcal S_{\mathcal K}$ is closed in $\GL(d)$, as
the preimage of the closed set $\mathcal K$ under $A \mapsto AA^{\tp}$; it contains $\Orth(d)$, as
$I \in \mathcal K$; and it is closed under products by \ref{prop:invariant-orders}(3). $\mathcal L$ is closed
under positive multiples and, $\mathcal S_{\mathcal K}$ being closed, under limits. It is closed under
sums by the Lie product formula $(e^{B/n}e^{B'/n})^n \to e^{B+B'}$. Put $P_n := e^{B/n}e^{B'/n}$,
$Q_n := e^{(B+B')/n}$ and $a := \|B\| + \|B'\|$. The curve $\varphi(s) := e^{sB}e^{sB'} - e^{s(B+B')}$ has
$\varphi(0) = 0$ and $\varphi'(0) = B + B' - (B + B') = 0$, so $n\|P_n - Q_n\| = n\|\varphi(1/n)\| \to 0$; and
$\|P_n\|, \|Q_n\| \le e^{a/n}$, as $\|e^Z\| \le e^{\|Z\|}$. From $P^n - Q^n = P(P^{n-1} - Q^{n-1}) + (P - Q)Q^{n-1}$,
by induction,
\[
  \|P_n^n - Q_n^n\| \le n\,(e^{a/n})^n\,\|P_n - Q_n\| = e^a\,n\|P_n - Q_n\| \to 0 ,
\]
and $Q_n^n = e^{B+B'}$. Applied to $tB$ and $tB'$, the products $(e^{tB/n}e^{tB'/n})^n$ lie in $\mathcal S_{\mathcal K}$ and tend to $e^{t(B+B')}$. Three cases then describe $\mathcal L$.
\begin{itemize}
\item For symmetric $X$, $e^{tX}(e^{tX})^{\tp} = e^{2tX}$, so $X \in \mathcal L$ exactly when
$X \in C(\mathcal K)$.
\item For skew $K$, $e^{tK} \in \Orth(d) \subseteq \mathcal S_{\mathcal K}$, so $\pm K \in \mathcal L$.
\item Adding the skew part of $B$ or its negative, $B \in \mathcal L$ if and only if
$\frac12(B + B^{\tp}) \in \mathcal L$.
\end{itemize} So $\mathcal L$ is as stated, and $C(\mathcal K) = \mathcal L \cap
\Sym(d)$ is a convex cone; it is closed, invariant and pointed by \ref{prop:invariant-orders}(4). By
(1), $\preceq_{C(\mathcal K)}$ is $\le_{\exp C(\mathcal K)}$, and $\exp C(\mathcal K) \subseteq \mathcal K$ by
definition. (c)$\Rightarrow$(b): let $M \in \mathcal K$ and $X := \log M$. Then $e^{tX} = M^t$ lies in
$\mathcal K$ for $t \in [0,1]$ by (c). For $t = j + \tau$ with $j \ge 1$ an integer and $0 \le \tau < 1$,
induction on $j$ gives it: $M^{j+\tau} = M^{1/2}M^{j-1+\tau}M^{1/2}$ is the composite of the steps $M$ and
$M^{j-1+\tau}$ of $\mathcal K$, which lies in $\mathcal K$ by transitivity
(\ref{prop:invariant-orders}(3)). So $X \in C(\mathcal K)$ and $\mathcal K \subseteq \exp C(\mathcal K)$. (b)$\Rightarrow$(a) is (1) with $C' = C(\mathcal K)$.
(a)$\Rightarrow$(c): by (1), $\mathcal K = \exp C'$, closed under $M \mapsto M^\tau$ as $C'$ is a cone, and
$C' = C(\exp C') = C(\mathcal K)$.

\emph{(3).} For $M, N$ with least eigenvalues at least $2$ and a unit $x$,
$\langle x, M^{1/2}NM^{1/2}x\rangle \ge 2|M^{1/2}x|^2 \ge 4$, so composites stay in $\mathcal K$, and
composites with $I$ trivially do: $\le_{\mathcal K}$ is transitive (\ref{prop:invariant-orders}(3)), and
reflexive. $\mathcal K$ is a step set, and closed since $\lambda^\uparrow_1$ is continuous. If $M \ne I$ and
$M \in \mathcal K$, then $\lambda^\uparrow_1(M^{-1}) = 1/\lambda^\uparrow_d(M) \le 1/2$, so $M^{-1} \notin
\mathcal K$: antisymmetry. For $X \ne 0$ and small $t > 0$, $e^{tX} \ne I$ and its least eigenvalue
$e^{t\lambda^\uparrow_1(X)}$ is below $2$, so $X \notin C(\mathcal K)$. A $\{0\}$-causal curve has
derivative $0$ almost everywhere, and an absolutely continuous curve with derivative $0$ almost everywhere is
constant, so $\preceq_{\{0\}}$ is equality. If $\le_{\mathcal K}$ were $\preceq_{C'}$, then for $X \in C'$ and
$t \ge 0$ the geodesic of (1) from $I$ to $e^{tX}$ would be $C'$-causal, so $I \le_{\mathcal K} e^{tX}$, that is
$e^{tX} \in \mathcal K$: $C' \subseteq C(\mathcal K) = \{0\}$. Then $\le_{\mathcal K}$ would be equality,
while $I \le_{\mathcal K} 2I$.

\emph{(4).} By \ref{lem:spectral-cones}(4) every nonzero $X \in C$ has $\tr X > 0$.
On $\PD(d)$ the second case is contained in the third, since $\Sigma' = c\Sigma$ has relative step
$cI = e^{(\log c)I}$ with $(\log c)I \in C$. If $\Sigma$ is definite, so is every $\Sigma'$ above it. If
$\Sigma \ne 0$ is singular, the $\Sigma'$ above it are its multiples $c\Sigma$, $c \ge 1$. The claims follow.
\begin{itemize}
\item Reflexivity is $c = 1$.
\item Transitivity: from $0$ trivially, from a nonzero singular matrix along its multiples, and on
$\PD(d)$ by that of $\preceq_C$.
\item Antisymmetry: on $\PD(d)$ it is (1). $\Sigma' \preceq^+_C 0$ forces $\Sigma' = 0$. And $c\Sigma = \Sigma$
with $\Sigma = c'(c\Sigma)$, $c, c' \ge 1$, forces $c = 1$.
\item Invariance: $A \cdot 0 = 0$, $A \cdot (c\Sigma) = c\,A\cdot\Sigma$, and $\preceq_C$ is invariant by (1).
$G$-invariance follows from $\GL(d)$-invariance.
\item The rays: $c\Sigma \preceq^+_C c'\Sigma$ for $0 \le c \le c'$, by the first case if $c = 0$ and the
second otherwise.
\item Chains: $s\Sigma'$ from $0$ and $(1 + s(c-1))\Sigma$ along a ray are chains by
\ref{thm:cascade-infinitely-divisible}(3)(b). On $\PD(d)$ the geodesic $\gamma$ of (1) is one, by
\ref{thm:cascade-infinitely-divisible}(3)(c) applied to $\mathcal K = \exp C$, a step set whose order is a partial
order by (1) and which is closed under $M \mapsto M^\tau$, $0 \le \tau \le 1$, as $C$ is a cone; and
$\le_{\exp C}$ is $\preceq_C$ by (1). So (P) holds.
\item The Loewner order: if $C \subseteq \PSD(d)$, then $e^X \ge I$ for $X \in C$, so every pair of
$\preceq^+_C$ is a Loewner pair. If $X \in C$ is not positive semidefinite, $I \preceq^+_C e^X$ while
$e^X - I$ is not positive semidefinite.
\item The ray order: for $C = \{cI : c \ge 0\}$, $\exp C = \{cI : c \ge 1\}$, so on $\PD(d)$ the third case
is the second, and $\preceq^+_C$ is the ray order.
\end{itemize}
\end{proof}

% CHANGED Q-0281: split from thm:cone-orders (ADR-0025; records/formalization/PRICING.md): clause (2),
% copied verbatim. Its proof needs the Lie product formula, which Mathlib lacks, where clause (1)
% needs majorization; the parent's text and numbering are unchanged.
% CHANGED Q-0281 wave 2, group F: proved (AffineHemigroup.ConeOrdersClosed, with the Lie product
% formula in AffineHemigroup.LieProduct); [A] on A9 through thm:cone-orders(1).
\begin{lemma}[closed invariant orders and their tangent wedges]
\label{lem:cone-orders-closed}
\uses{thm:cone-orders,prop:invariant-orders,thm:cited-A9}
\lean{AffineHemigroup.closed_order_tangent_wedge,
AffineHemigroup.closed_order_cone_tfae}\leanok
Let $d \ge 1$, and use the notation of \ref{prop:invariant-orders} and \ref{thm:cone-orders}. Let
$\mathcal K$ be a closed step set for which $\le_{\mathcal K}$ is a partial order. Then
$\{B \in \gl(d) : e^{tB} \in \mathcal S_{\mathcal K} \text{ for every } t \ge 0\} =
\{B : \tfrac12(B + B^{\tp}) \in C(\mathcal K)\}$, and this is a closed convex cone containing every
skew-symmetric matrix; $C(\mathcal K)$ is an $\Orth(d)$-invariant closed convex cone with
$C(\mathcal K) \cap (-C(\mathcal K)) = \{0\}$; and $\preceq_{C(\mathcal K)}$ is contained in
$\le_{\mathcal K}$. The following are equivalent: (a) $\le_{\mathcal K}$ is $\preceq_{C'}$ for some cone
$C'$ as in \ref{thm:cone-orders}; (b) $\mathcal K = \exp C(\mathcal K)$; (c) $M^\tau \in \mathcal K$
whenever $M \in \mathcal K$ and $0 \le \tau \le 1$. When they hold, $C' = C(\mathcal K)$.
\statusA\quad\ledger{A9}\quad\emph{(\textbf{Assignment.} \ledger{A9} enters through
\ref{thm:cone-orders}(1), which gives the inclusion of $\preceq_{C(\mathcal K)}$ in $\le_{\mathcal K}$ and
the equivalence of (a) with (b) and (c). The tangent wedge, its description and the cone $C(\mathcal K)$
are $[T]$, with the Lie product formula proved in the parent's proof. Clause (2) of
\ref{thm:cone-orders}, split off on 2026-10-04 (Q-0281, ADR-0025); the numbering stays the parent's.)}
\end{lemma}

\begin{proof}
\leanok
This is the part of the proof of \ref{thm:cone-orders} headed (2).
\end{proof}

% CHANGED Q-0281: split from thm:cone-orders (ADR-0025; records/formalization/PRICING.md): clauses
% (3) and (4), copied with the theorem's standing hypothesis on C restated in (2) here and its clause
% (2) cited as lem:cone-orders-closed. A counterexample and the extension to PSD(d), of a different
% kind and cost from clauses (1) and (2); the parent's text and numbering are unchanged.
% CHANGED Q-0281: \uses re-pointed to the split child it reads (thm:cascade-infinitely-divisible (3)(b)-(c), in the parent's proof of (4), beside (2)).
% CHANGED Q-0281 wave 2: blind review W2F-5: clause (1)'s parenthetical names the geodesic of (1) of
% thm:cone-orders, as in the parent, and the \uses edge to lem:cone-orders-closed is dropped: nothing
% else of this node reads it (the proof of (3) uses the geodesic, that of (4) reads (1) only).
% CHANGED Q-0281 wave 2, group F: proved (AffineHemigroup.ConeOrdersPSD); [A] on A9 through
% thm:cone-orders(1) for clause (2); clause (1) Lean core.
% CHANGED Q-0469: clause (1)'s parenthetical goes, as in the parent's clause (3) (a proof inside the
% statement; records/affine/GATE-affine.md, Q-0317 row "§4 22"); the parent's proof of (3) carries it.
\begin{lemma}[a closed order without a cone; the orders on $\PSD(d)$]
\label{lem:cone-orders-psd}
\uses{thm:cone-orders,lem:spectral-cones,def:gw-cascade-family,thm:cascade-infinitely-divisible,lem:path-condition-orders,thm:cited-A9}
\lean{AffineHemigroup.closed_order_without_cone, AffineHemigroup.causalPlus_wedge,
AffineHemigroup.causalPlus_ray}\leanok
Let $d \ge 1$, and use the notation of \ref{prop:invariant-orders}, \ref{lem:spectral-cones} and
\ref{thm:cone-orders}.
\begin{enumerate}
\item \emph{(A closed order without a cone.)} $\mathcal K := \{I\} \cup \{M \in \PD(d) :
\lambda^\uparrow_1(M) \ge 2\}$ is a closed step set, $\le_{\mathcal K}$ is a partial order, and
$C(\mathcal K) = \{0\}$, so that $\preceq_{C(\mathcal K)}$ is equality; hence $\le_{\mathcal K}$ is
$\preceq_{C'}$ for no cone $C'$.
\item \emph{(On $\PSD(d)$.)} Let $C \subseteq \Sym(d)$ be a closed convex cone invariant under
$X \mapsto QXQ^{\tp}$, $Q \in \Orth(d)$, and suppose $I \in C$ and $C \cap (-C) = \{0\}$. For
$\Sigma, \Sigma' \in \PSD(d)$ let $\Sigma \preceq^+_C \Sigma'$ when $\Sigma = 0$, or $\Sigma' = c\,\Sigma$
for some $c \ge 1$, or $\Sigma, \Sigma' \in \PD(d)$ and $\Sigma \preceq_C \Sigma'$. Then $\preceq^+_C$ is
a $\GL(d)$-invariant partial order on $\PSD(d)$ that contains the rays, every pair
$\Sigma \preceq^+_C \Sigma'$ is joined by a chain as in \ref{thm:cascade-infinitely-divisible}(2), and
$(G, \preceq^+_C)$ is a wedge datum with the path condition (P) for every closed subgroup $G$ of
$\GL(d)$ containing the dilations. It is contained in the Loewner order if and only if
$C \subseteq \PSD(d)$. For $C = \{cI : c \ge 0\}$ it is the ray order of \ref{def:gw-cascade-family}.
\end{enumerate}
\statusA\quad\ledger{A9}\quad\emph{(\textbf{Assignment.} \ledger{A9} enters clause (2) through
\ref{thm:cone-orders}(1), the order $\preceq_C$ on $\PD(d)$; the last sentence of (2), the ray order, reads (1)
the same way. Clause (1) is $[T]$ and spends nothing: it reads only the geodesic of (1). Clauses (3) and (4)
of \ref{thm:cone-orders}, split off on 2026-10-04 (Q-0281, ADR-0025); the numbering stays the parent's.)}
\end{lemma}

\begin{proof}
\leanok
This is the part of the proof of \ref{thm:cone-orders} headed (3) and (4).
\end{proof}

% not in the draft: the plan's item D3, its named question (Q-0266): is the stable step cone the order
% generated by the cone field K_alpha? Yes, exactly; one inclusion was noted in Q-0265 (the geodesic of
% thm:cascade-infinitely-divisible(3)(c)), the other is thm:cone-orders(1). What writing it down found:
% the stable order is the Loewner order read through the map Sigma -> (det Sigma)^{-1/(d+alpha)} Sigma
% (clause (3)), which gives a second, one-line proof of its transitivity and of its generation, with no
% Levy measure.
% CHANGED Q-0281: \uses re-pointed to the split child it reads (condition (c) and the order (2) of prop:stable-step-cone only; thm:cone-orders(2) beside (1)).
% CHANGED Q-0281 wave 2, group F: proved (AffineHemigroup.StableOrderGenerated) by the entry-free
% route, the author's decision of 2026-10-05: the proof is rewritten so that it reads neither
% thm:cone-orders(1)'s converse (which spends ledger A9) nor lem:cone-orders-closed, and \uses drops
% both (the second is the session's narrowing to what the proof reads). The order of a cone field is
% still the one thm:cone-orders defines; the statement is unchanged. The Lean prints Lean core only.
% CHANGED Q-0327: wording and notation only (whole-paper review flags 17-20; blind draft review flags
% 5, 27, 28): the title names the cone field $C_\alpha$ (was $K_\alpha$); "Three orders" for "Three
% corners"; "tangent cone", defined in (2), for "tangent wedge"; the order of a cone field is
% $\preceq_{C_\alpha}$ throughout, so that $\preceq_\alpha$ is defined once, by
% prop:stable-step-cone(2), and (1) shows that the two agree for $0 < \alpha < 2$.
% CHANGED round-2 review (Gemini), item 12, the author's decision D7 (2026-11-03): the exponent cone K_alpha
% is written frak k_alpha (a subset of gl(d), in fraktur as gl(d) is), apart from the step cone
% cal K_alpha. Notation only; no claim, \uses or \lean changes.
\begin{proposition}[the stable step cone is the order of the cone field $C_\alpha$]
\label{prop:stable-order-generated}
\uses{lem:stable-step-order,prop:stable-exponent-cone}
\lean{AffineHemigroup.stableGenCone_props, AffineHemigroup.psdSet_cone_props,
AffineHemigroup.stable_order_zero, AffineHemigroup.stable_order_pos,
AffineHemigroup.stable_order_infty, AffineHemigroup.stableExpCone_tangent,
AffineHemigroup.stable_order_loewner_coords}\leanok
Let $d \ge 1$. For $0 \le \alpha < \infty$ put
\[
  C_\alpha := \{X \in \Sym(d) : X \text{ positive semidefinite},\ (d+\alpha)\,\lambda_{\min}(X) \ge \tr X\} ,
  \qquad C_\infty := \PSD(d) ,
\]
and write $\preceq_{C_\alpha}$ for the order of the cone field $C_\alpha$ (\ref{thm:cone-orders}), $\lambda_{\min}$ the least
eigenvalue. Each $C_\alpha$ is an $\Orth(d)$-invariant closed convex cone containing $I$, with
$C_\alpha \cap (-C_\alpha) = \{0\}$. For $\Sigma, \Sigma' \in \PD(d)$ let
$M := \Sigma^{-1/2}\Sigma'\Sigma^{-1/2}$.
\begin{enumerate}
\item \emph{(Three orders, each generated by its cone field.)} $C_0 = \{cI : c \ge 0\}$, and
$\Sigma \preceq_{C_0} \Sigma'$ if and only if $\Sigma' = c\,\Sigma$ with $c \ge 1$, the ray order on $\PD(d)$.
For $0 < \alpha < \infty$ the condition $X \ge 0$ in $C_\alpha$ is implied by the inequality, and
$\Sigma \preceq_{C_\alpha} \Sigma'$ if and only if $\lambda_{\min}(M)^{d+\alpha} \ge \det M$; for
$0 < \alpha < 2$ this is the order $\preceq_\alpha$ of \ref{prop:stable-step-cone}(2). $\Sigma \preceq_{C_\infty}
\Sigma'$ if and only if $\Sigma \le \Sigma'$, the Loewner order. In each case the related pairs are exactly
the pairs joined by a $C_\alpha$-causal curve, and $\Sigma^{1/2}M^t\Sigma^{1/2}$ is one.
\item \emph{(The exponent cone is the tangent cone.)} For $0 < \alpha < 2$ the set $\mathfrak k_\alpha$ of
\ref{eq:stable-exponent-cone} equals $\{B \in \gl(d) : e^{tB}e^{tB^{\tp}} \in \exp C_\alpha \text{ for every }
t \ge 0\}$, the tangent cone of the semigroup $\{A : I \preceq_{C_\alpha} AA^{\tp}\}$, the set of $B$ with $e^{tB}$ in
it for every $t \ge 0$, and
$\mathfrak k_\alpha \cap \Sym(d) = C_\alpha$.
\item \emph{(The Loewner order in other coordinates.)} For $0 < \alpha < \infty$ let
$\Phi_\alpha(\Sigma) := (\det\Sigma)^{-1/(d+\alpha)}\,\Sigma$ on $\PD(d)$. Then $\Phi_\alpha$ is a
diffeomorphism of $\PD(d)$ onto itself with inverse $T \mapsto (\det T)^{1/\alpha}\,T$;
$\Sigma \preceq_{C_\alpha} \Sigma'$ if and only if $\Phi_\alpha(\Sigma) \le \Phi_\alpha(\Sigma')$; the derivative of
$\Phi_\alpha$ at $\Sigma$ maps the cone $\Sigma^{1/2}C_\alpha\Sigma^{1/2}$ onto $\PSD(d)$; and
$\Phi_\alpha(A\Sigma A^{\tp}) = \psi_\alpha(A)\,\Phi_\alpha(\Sigma)\,\psi_\alpha(A)^{\tp}$ for
$A \in \GL(d)$, where $\psi_\alpha(A) := |\det A|^{-1/(d+\alpha)}A$ is an automorphism of the group
$\GL(d)$. In particular $\psi_\alpha$ maps the semigroup $\{A : I \preceq_{C_\alpha} AA^{\tp}\}$ onto the semigroup
$\{A : AA^{\tp} \ge I\}$ of the Loewner order.
\end{enumerate}
\statusT\quad\emph{(Cites nothing, and spends no ledger entry: the proof reads the stable orders
through $\Phi_\alpha$ and the Loewner order, not through \ref{thm:cone-orders}(1), whose superadditivity
spends A9 (the author's decision of 2026-10-05). The set
$\mathfrak k_\alpha$ is used here as the set defined by \ref{eq:stable-exponent-cone}; that it is the exponent set of
$S_\alpha(I)$ is \ref{prop:stable-exponent-cone}'s content and spends \ledger{A3} there, not here.
\textbf{What it is worth.} It answers the plan's item D3 as posed: the finite stable step cone
$\lambda_{\min}(M)^{d+\alpha} \ge \det M$ of \ref{prop:stable-step-cone} is exactly the order generated by the
cone field of $\mathfrak k_\alpha \cap \Sym(d)$, with no closure, and the same holds for the rays and for the Loewner
order. One inclusion was noted with \ref{thm:cascade-infinitely-divisible}(3)(c), whose chain is the
geodesic of (1); the other is the monotonicity of $\Phi_\alpha$ along a causal curve, which
\ref{thm:cone-orders}(1) also gives. Clause (2) says in what sense $\mathfrak k_\alpha$ is the
infinitesimal cone: it is the tangent wedge of the semigroup of the order, with the skew matrices as its
edge, which is why \ref{prop:stable-exponent-cone} found every skew matrix in it. Clause (3) shows that the
stable orders are not new orders up to a change of coordinates: $\preceq_{C_\alpha}$ is the Loewner order read
through $\Phi_\alpha$, and the semigroup of $\preceq_{C_\alpha}$ is the preimage under the automorphism
$\psi_\alpha$ of the semigroup of the Loewner order, which is the set of inverses of the contraction semigroup
of \sscite{hilgert1989lie} Example~V.4.55, pp.~437--438 (read, not used). This gives a second proof of the transitivity of $\preceq_\alpha$, which
\ref{prop:stable-step-cone}(2) obtained from the ordering of \Levy{} densities. The parameter interpolates
between the rays ($\alpha \to 0$, where $\Phi_\alpha$ degenerates to the shape map $\Sigma \mapsto
(\det\Sigma)^{-1/d}\Sigma$) and the Loewner order ($\alpha = \infty$). \textbf{Not claimed}: that $\Phi_\alpha$
is related to any probabilistic operation on the kernels.)}
\end{proposition}

% CHANGED Q-0281 wave 2, group F: proof rewritten to the entry-free route (the author's decision of
% 2026-10-05): the two facts about curves, Phi_alpha along a causal curve, and the geodesic, in place
% of thm:cone-orders(1)-(2).
\begin{proof}
\leanok
\emph{The cones.} The conditions defining $C_\alpha$ are invariant under orthogonal conjugation and positive
scaling; $\lambda_{\min}$ is concave and continuous and $\tr$ is linear, so $C_\alpha$ is a closed convex
cone; $I$ satisfies $d + \alpha \ge d$; and $C_\alpha \subseteq \PSD(d)$ meets $-\PSD(d)$ only in $0$.

\emph{Three facts about curves.} (a) Let $\eta : [0,1] \to \Sym(d)$ have absolutely continuous entries. If
$\eta'(t)$ is positive semidefinite for almost every $t$, then for each $x$ the function $\langle x,
\eta(t)x\rangle$ is absolutely continuous with derivative $\langle x, \eta'(t)x\rangle \ge 0$ almost
everywhere, so it is the integral of its derivative and $\eta(0) \le \eta(1)$. If $\eta' = 0$ almost
everywhere, each entry is constant, and $\eta(0) = \eta(1)$. (b) For $s \in \RR$ and a $C$-causal $\gamma$
($C$ any cone), the entries of $(\det\gamma)^s\gamma$ are absolutely continuous on $[0,1]$: $\det\gamma$ is a
polynomial in the entries of $\gamma$, positive and continuous, hence between two positive numbers on
$[0,1]$, where $x \mapsto x^s$ is Lipschitz. At almost every $t$ the entries of $\gamma$ are differentiable;
there, writing $\gamma(u) = \gamma(t)\cdot\gamma(t)^{-1}\gamma(u)$ and expanding the determinant of
$\gamma(t)^{-1}\gamma(u)$, which is $I$ at $u = t$, gives Jacobi's formula $(\det\gamma)'(t) =
\det\gamma(t)\,\tr(\gamma(t)^{-1}\gamma'(t))$. With $Y := \gamma^{-1/2}\gamma'\gamma^{-1/2}$, so that
$\gamma' = \gamma^{1/2}Y\gamma^{1/2}$ and $\tr(\gamma^{-1}\gamma') = \tr Y$,
\[
  \bigl((\det\gamma)^s\gamma\bigr)' = (\det\gamma)^s\bigl(\gamma' + s(\tr Y)\gamma\bigr)
  = (\det\gamma)^s\,\gamma^{1/2}\bigl(Y + s(\tr Y)I\bigr)\gamma^{1/2} .
\]
(c) If $C$ is invariant under orthogonal conjugation, $X \in C$ and $M = e^X$ is the relative step of
$(\Sigma, \Sigma')$, then $\gamma(t) := \Sigma^{1/2}M^t\Sigma^{1/2} = \Sigma^{1/2}e^{tX}\Sigma^{1/2}$ is a
$C$-causal curve from $\Sigma$ to $\Sigma'$: its entries are smooth, and with $A := \Sigma^{1/2}e^{tX/2}$,
$\gamma(t) = AA^{\tp}$ and $\gamma'(t) = AXA^{\tp}$, so $\gamma(t)^{-1/2}\gamma'(t)\gamma(t)^{-1/2} =
QXQ^{\tp}$ with $Q := \gamma(t)^{-1/2}A$ orthogonal.

\emph{(1).} $X \in C_0$ means $X \ge 0$ and $d\lambda_{\min}(X) \ge \tr X$, that is, every eigenvalue equals the
least one and is nonnegative: $X = cI$, $c \ge 0$. Let $\gamma$ be $C_0$-causal from $\Sigma$ to $\Sigma'$.
Almost everywhere $Y = cI$ for some $c \ge 0$, so $Y + s(\tr Y)I = 0$ for $s = -1/d$, and by (b) and (a)
$(\det\gamma)^{-1/d}\gamma$ is constant: $\Sigma' = c\,\Sigma$ with $c = (\det\Sigma'/\det\Sigma)^{1/d}$. As
$C_0 \subseteq \PSD(d)$, $\gamma' = \gamma^{1/2}Y\gamma^{1/2}$ is positive semidefinite almost everywhere, so
$\Sigma \le \Sigma'$ by (a), and $(c - 1)\Sigma \ge 0$ forces $c \ge 1$. Conversely the relative step of
$(\Sigma, c\Sigma)$ is $cI = e^{(\log c)I}$ with $(\log c)I \in C_0$ for $c \ge 1$, and (c) gives the curve.
For $0 < \alpha < \infty$, $\tr X \ge d\lambda_{\min}(X)$, so
$(d+\alpha)\lambda_{\min}(X) \ge \tr X$ gives $\alpha\lambda_{\min}(X) \ge 0$, hence $X \ge 0$. With
$M = e^X$, $\lambda_{\min}(M) = e^{\lambda_{\min}(X)}$ and $\det M = e^{\tr X}$, so $M \in \exp C_\alpha$ if and
only if $\lambda_{\min}(M)^{d+\alpha} \ge \det M$, which is condition (c) of \ref{prop:stable-step-cone}(1)
when $0 < \alpha < 2$. Let $\gamma$ be $C_\alpha$-causal and $s := -1/(d+\alpha)$. Almost everywhere $Y \in
C_\alpha$, so $Y + s(\tr Y)I = Y - \frac{\tr Y}{d+\alpha}I \ge 0$, and by (b) and (a) $\Phi_\alpha(\Sigma) \le
\Phi_\alpha(\Sigma')$. Dividing by $(\det\Sigma)^s > 0$ and congruence by $\Sigma^{-1/2}$, $\Phi_\alpha(\Sigma)
\le \Phi_\alpha(\Sigma')$ is $I \le (\det M)^sM$, that is $\lambda_{\min}(M)(\det M)^s \ge 1$, which is
$\lambda_{\min}(M)^{d+\alpha} \ge \det M$. Conversely, if that holds, $M \in \exp C_\alpha$ and (c) gives the
curve. Finally $e^X \ge I$ if and only if $X \ge 0$, so $\exp C_\infty = \{M : M \ge I\}$, and
$M \ge I$ is $\Sigma \le \Sigma'$, congruence preserving positive semidefiniteness. A $\PSD(d)$-causal curve has
$\gamma' = \gamma^{1/2}Y\gamma^{1/2} \ge 0$ almost everywhere, so $\Sigma \le \Sigma'$ by (a); conversely
$\Sigma \le \Sigma'$ puts $M$ in $\exp C_\infty$, and (c) gives the curve.

\emph{(2).} Let $s := -1/(d+\alpha)$ and, for a matrix $B$, $B' := B + s(\tr B)I$. As
$|\det e^{tB}|^s = e^{st\tr B}$ and $I$ commutes with $B$, $\psi_\alpha(e^{tB}) = e^{tB'}$. By (1) and (3),
$e^{tB}e^{tB^{\tp}} \in \exp C_\alpha$ if and only if $I \preceq_{C_\alpha} e^{tB}(e^{tB})^{\tp}$, if and only if
$I \le \Phi_\alpha(e^{tB}(e^{tB})^{\tp}) = \psi_\alpha(e^{tB})\psi_\alpha(e^{tB})^{\tp} = e^{tB'}e^{tB'^{\tp}}$.
For $x \in \RR^d$ put $f(t) := |e^{tB'^{\tp}}x|^2 = \langle x, e^{tB'}e^{tB'^{\tp}}x\rangle$, whose derivative
is $\langle y, (B' + B'^{\tp})y\rangle$ with $y = e^{tB'^{\tp}}x$. If $\lambda_{\min}(\frac12(B' + B'^{\tp}))
\ge 0$ then $f(t) \ge f(0) = |x|^2$ for $t \ge 0$; conversely $f(t) \ge f(0)$ for $t \ge 0$ gives
$f'(0) = \langle x, (B' + B'^{\tp})x\rangle \ge 0$. So $e^{tB}e^{tB^{\tp}} \in \exp C_\alpha$ for every
$t \ge 0$ if and only if $\lambda_{\min}(\frac12(B + B^{\tp})) + s\tr B \ge 0$, that is
$(d+\alpha)\lambda_{\min}(\frac12(B + B^{\tp})) \ge \tr B$, which is \ref{eq:stable-exponent-cone}. A symmetric
$X$ has $\frac12(X + X^{\tp}) = X$, and the condition $X \ge 0$ is implied for $\alpha > 0$, so
$\mathfrak k_\alpha \cap \Sym(d) = C_\alpha$.

\emph{(3).} Put $s := -1/(d+\alpha)$, so $1 + ds = \alpha/(d+\alpha) > 0$. For $T := \Phi_\alpha(\Sigma)$,
$\det T = (\det\Sigma)^{1+ds}$, so $(\det T)^{1/\alpha} = (\det\Sigma)^{1/(d+\alpha)}$ and
$(\det T)^{1/\alpha}T = \Sigma$; in the same way $\Phi_\alpha((\det T)^{1/\alpha}T) = T$. Both maps are
smooth. By (1), $\Sigma \preceq_{C_\alpha} \Sigma'$ if and only if $\Phi_\alpha(\Sigma) \le \Phi_\alpha(\Sigma')$.
By Jacobi's formula of (b) applied to $\varepsilon \mapsto \Sigma + \varepsilon V$, the derivative of
$\Phi_\alpha$ at $\Sigma$ in the direction $V$ is
\[
  D\Phi_\alpha(\Sigma)V = (\det\Sigma)^s\bigl(V + s\tr(\Sigma^{-1}V)\,\Sigma\bigr) ,
  \qquad
  D\Phi_\alpha(\Sigma)\bigl(\Sigma^{1/2}Y\Sigma^{1/2}\bigr) = (\det\Sigma)^s\,\Sigma^{1/2}\bigl(Y + s(\tr Y)I\bigr)\Sigma^{1/2} .
\]
The second is positive semidefinite when $\lambda_{\min}(Y) \ge (\tr Y)/(d+\alpha)$, in particular for
$Y \in C_\alpha$. Conversely, for $P \ge 0$ put $Z := (\det\Sigma)^{-s}\Sigma^{-1/2}P\Sigma^{-1/2} \ge 0$ and
$Y := Z - \frac{s}{1+ds}(\tr Z)I$; then $\tr Y = \tr Z/(1+ds)$ and $Y + s(\tr Y)I = Z \ge 0$, so $Y \in
C_\alpha$, and $D\Phi_\alpha(\Sigma)(\Sigma^{1/2}Y\Sigma^{1/2}) = P$: every positive semidefinite direction is
attained.
For $A \in \GL(d)$, $\det(A\Sigma A^{\tp}) = (\det A)^2\det\Sigma$ gives the covariance. Finally
$\psi_\alpha(AB) = \psi_\alpha(A)\psi_\alpha(B)$, $|\det A|^s$ being a character and scalars central, and with
$s' := -s/(1+ds)$ the map $A \mapsto |\det A|^{s'}A$ is its inverse, since $|\det(|\det A|^{s'}A)|^s|\det A|^{s'} =
|\det A|^{s'(1+ds) + s} = 1$, and symmetrically for the other composition. As $\Phi_\alpha(I) = I$, $I \preceq_{C_\alpha} AA^{\tp}$ if and only if
$I \le \Phi_\alpha(AA^{\tp}) = \psi_\alpha(A)\psi_\alpha(A)^{\tp}$, which is the last sentence.
\end{proof}
